% ============================================================================
\clearpage
\section*{v4.9.0-r26 Source-Only: Pre-Registered Test of Degree-Six Potentials on the Triple Space}
\addcontentsline{toc}{section}{v4.9.0-r26 Source-Only: Pre-Registered Test of Degree-Six Potentials on the Triple Space}
\label{sec:v490-r26-degree-six-triple-space}
% ============================================================================

\paragraph{Purpose.}
Layer v4.8.0-r16 showed that invariant potentials up to degree four on the triple space of the triality slots give at
most two distinct slot norms, and named degree six as necessary for three.  This layer records the pre-registered test
H16 of degree six (note \texttt{llm/Vorab\_Notiz\_H16\_Grad6\_Tripelraum.md}, commit \texttt{d0360b8}, fixed before
computing; created under the label H12, which was already taken, and renamed without change of content in commit
\texttt{6bd2f1b}).  Earlier layers are not rewritten.

\noindent\textbf{Source-only status.}
This v4.9.0-r26 layer is source-only.  It does not build a PDF, does not assign a new DOI,
and does not authorize a public release.  The v4.9.0 Version DOI
\href{https://doi.org/10.5281/zenodo.23534954}{10.5281/zenodo.23534954} and the Concept DOI
\href{https://doi.org/10.5281/zenodo.19210728}{10.5281/zenodo.19210728} are unchanged by it.

\subsection*{v4.9.0-r26 setting and hypotheses}
\label{subsec:v490-r26-setting}

As in v4.8.0-r16: slot norms \(n_k\), \(x_k=n_k^2\), \(S=\sum x_k\), \(Q=\sum x_k^2\), \(H=\sum x_k^3\),
\(P=n_1n_2n_3\), \(t=\operatorname{Re}((ab)c)=P\) at the minimum, reading \(\sqrt{m_k}\propto n_k\).  Only \(H\) adds a
term \(x_k^3\) to the condition at critical points, so the fixed minimal extension is
\[
V_6=-S+\beta S^2+\gamma Q-\kappa P+\delta H,\qquad \delta>0,\ \kappa\ge0 .
\]
H16a: for part of a fixed parameter distribution (\(N=2000\)), \(V_6\) has a global minimum with three distinct non-zero
norms.  H16b: for some parameters the charged-lepton triple is the global minimum.  H16c: \(K=2/3\) is singled out among
such minima (at least half within \(0.01\)).

\subsection*{v4.9.0-r26 results}
\label{subsec:v490-r26-results}

Script \texttt{llm/k64\_test\_h16.py}.  Global minima over the distribution: 1585 with equal norms, 398 with one occupied
slot, 17 with a degenerate pair, \textbf{none} with three distinct norms.  The lepton triple is a critical point on the
whole two-parameter family given by Vieta's relations, but on the \(41\times41\) grid it is nowhere even a local minimum
(1139 points with \(\delta\le0\), 542 with an indefinite Hessian).

\paragraph{Verdict.}
H16a, H16b and H16c are refuted under the fixed criterion.

\paragraph{Theoremlet v4.9.0-r26.1 (found after the verdict).}
Every polynomial potential in the slot norms and \(P\) of degree at most six has the form
\(V=A(S,P)+e_2\,(b_0+b_1S)\) with \(e_2=\sum_{i<j}x_ix_j\).  At a critical point with three distinct non-zero norms it is
a saddle if \(b_1\neq0\); if \(b_0=b_1=0\) it does not depend on \(e_2\) and the norms are not fixed.  Hence no potential of
degree at most six has a minimum with three distinct norms; degree at least eight is necessary.

\paragraph{Proof sketch.}
\(Q=S^2-2e_2\), \(H=S^3-3Se_2+3P^2\), \(SQ=S^3-2Se_2\); the other invariants \(S,S^2,S^3,P,SP,P^2\) do not contain
\(e_2\).  For distinct \(x_k\) the map \(n\mapsto(S,e_2,P)\) is a local diffeomorphism (Vandermonde), and at a critical
point the Hessian transforms by congruence.  In these coordinates it reads
\(\left(\begin{smallmatrix}\ast&b_1&\ast\\ b_1&0&0\\ \ast&0&\ast\end{smallmatrix}\right)\), whose leading
\(2\times2\) minor is \(-b_1^2<0\).  For \(V_6\), \(b_1=-3\delta\), and \(\partial V/\partial e_2=0\) is the Vieta relation
\(2\gamma+3\delta S=0\).  Checked exactly with 40 digits at 2000 random critical points of the general degree-six ansatz
(script \texttt{llm/k65\_sattel\_grad6.py}); in double precision the negative eigenvalue can fall to \(5\cdot10^{-15}\) of
the largest.  \(\square\)

\paragraph{Reclassification.}
The statement of v4.8.0-r16 that degree six is necessary holds for critical points; for minima degree six is not enough.
Whether degree eight suffices is not shown.  It adds at least four coefficients, so the two lepton ratios could at best be
fitted, and \(K=2/3\) involves \(\sum n_k\), which is not a polynomial invariant.  Within this framework the generation
\(S_3\) with natural invariants does not produce a mass hierarchy; one needs an explicit breaking of \(S_3\) or invariants
beyond the norms, such as the real parts of the individual slots, which \(G_2\) leaves fixed.

\subsection*{v4.9.0-r26 status in the three frames}
\label{subsec:v490-r26-frames}

\begin{center}
\small
\begin{tabular}{@{}>{\raggedright\arraybackslash}p{0.30\textwidth}>{\raggedright\arraybackslash}p{0.62\textwidth}@{}}
\toprule
result & M / T / P\\
\midrule
H16 verdict, r26.1 & frame-free: statements about octonion norms and polynomial invariants, no frame quantity enters\\
mass reading \(\sqrt{m_k}\propto n_k\) & matter frame M (charged-lepton masses, as the Koide phase in v4.9.0-r22); not evaluated in T or P\\
\bottomrule
\end{tabular}
\end{center}

\subsection*{v4.9.0-r26 ledger and guard update}
\label{subsec:v490-r26-ledger-guard}

\[
R_{\rm cg}^{v490\text{-}r26}=(0,0,1)
\]
for (separation consistent with \(G_2\); carrier of the triality slots identified; mass mechanism open, with all
norm-invariant potentials up to degree six excluded).

The layer spends no guard:
\[
R_{\rm glob}^{(v490\text{-}r26,+)}=(0,1,1,1,1),
\qquad
\Gamma_{\rm guard}^{v490\text{-}r26}=(1,1,1,1,1),
\qquad
\beta_4=1/256\ \text{unspent}.
\]
It does not imply exact finite \(C5\) zero closure, exact all-readout terminality,
\(C6/T2\), a global photon-clock/RH theorem, or global all-readout closure.
